% texts/openstax-calculus1.tex -- THE TEXT LAYER FOR OPENSTAX CALCULUS VOLUME 1.
%
% Everything here answers one question: WHAT DOES OPENSTAX CALCULUS VOLUME 1
% NAME THINGS, AND WHERE DOES IT LIVE? Nothing else belongs in this file.
%
% THE FILE IS NAMED FOR THE BOOK, NOT THE PUBLISHER. It was openstax.tex while
% it lived in one course. In the shared machinery a name must be unique across
% every course, and OpenStax publishes many books -- Algebra and Trigonometry
% among them, which this same course cites. THE BOOK'S KEY IS STILL `openstax':
% \booktitle{openstax}, \SetDefaultText{openstax} below and the \DeclareBook in
% books.tex are unchanged. Only the argument of \usetext follows the file name.
%
% WRITTEN FOR MATH 1003, and the comments below still speak from there: "this
% course" and its modules are MATH 1003's. There, this is the text for MODULE A
% ONLY -- the five review lectures on functions (TOPICS.tsv rows A.1-A.5, src
% O1.1-O1.5). Modules B, C and D are APEX and use apex.tex beside this one.
%
% WHY THERE ARE TWO FILES HERE. The noun for a worked item and the result
% vocabulary below are DOCUMENT-WIDE settings, so a document has to pick one
% book; the two files cannot be merged. \usetext is that pick, and it is why
% every Module A master carries
%
%   \usetext{openstax-calculus1}
%
% There is no course default to fall back on: textbook.tex was deleted when this
% file was split out, precisely so that naming the book is mandatory. A master
% that names nothing gets no vocabulary and no base, and says so.
%
% Because only one of these files is ever loaded into a document, the two may
% freely disagree -- different nouns, and result environments of the same or
% different names -- with no risk of collision. The corollary bites in one place
% and is easy to forget: a Module A lecture never loads apex.tex, so APEX's
% cross-reference list is not available to it, and vice versa. See §4.


%%%%%%%%%%%%%%%%%%%%%%%%%%%%%%%%%%%%%%%%
% 1. WHAT THE TEXT NAMES ITS RESULTS
%%%%%%%%%%%%%%%%%%%%%%%%%%%%%%%%%%%%%%%%
% TWO ARE DECLARED, at the foot of this section: Rule and Strategy, each added
% when a Module A lecture first needed it. The Checkpoint is still NOT declared,
% and that is a decision rather than an omission.
%
% Deliberately NOT "Key Idea": that is APEX's word, declared in apex.tex, and a
% document built from this file has no \keyidea at all. OpenStax's own flagged,
% numbered result is the Checkpoint -- the prompt that follows each worked
% Example -- and the line to add, if a lecture ever wants one, is:
%
%   \NewTextResult[cbdefn]{checkpoint}{Checkpoint}{cb/textresult}
%
% (cbdefn, upright body, rather than cbthm's italic: a checkpoint is a prompt to
% stop and try something, not a stated result.)
%
% It stays commented because no Module A lecture uses one. The old ch0 sources
% reach for \begin{definition} and for undifferentiated coloured boxes, and the
% coloured ones are being read case by case during conversion -- most land on
% CourseBoxes' own vocabulary (definition, definition*, remark*) or on an aside
% (takeaway, tip). Declare a Checkpoint when a lecture actually needs one, on
% the same "declare only what you cite" discipline as §4.
%
% GUARDED, if it is ever uncommented: \NewTextResult comes from CourseBoxes.sty,
% which Assessment.cls does not load. The guard lets this file be \input by ANY
% master.
%
\ifdefined\NewTextResult
  \NewTextResult[cbdefn]{Rule}{Rule}{cb/textresult}
\fi
\ifdefined\NewTextResult
  \NewTextResult[cbdefn]{strategy}{Strategy}{cb/textresult}
\fi


%%%%%%%%%%%%%%%%%%%%%%%%%%%%%%%%%%%%%%%%
% 2. WHAT THE TEXT CALLS A WORKED ITEM
%%%%%%%%%%%%%%%%%%%%%%%%%%%%%%%%%%%%%%%%
% OpenStax says Example, which is already the machinery's default -- the same
% answer APEX gives, arrived at independently. So this stays commented, here as
% the documented place to change it rather than as a setting that needs making.
% \renewcommand{\NotesProblemLabelWord}{Exercise}


%%%%%%%%%%%%%%%%%%%%%%%%%%%%%%%%%%%%%%%%
% 3. WHICH BOOK THIS DOCUMENT IS BUILT ON
%%%%%%%%%%%%%%%%%%%%%%%%%%%%%%%%%%%%%%%%
% One line, and it is a statement about the DOCUMENT, not about the book: every
% \booksection and \textbookref that names no book resolves against this one.
%
% THE BASE URL IS NOT HERE ANY MORE -- it moved to books.tex, which every
% document loads whole. Where OpenStax lives is a fact about OpenStax; only the
% pick of vocabulary is exclusive to a document, and this line is that pick.
%
% (The observation that travelled to books.tex with the URL: OpenStax page slugs
% carry their own "pages/" segment and NO .html extension, unlike APEX's. Slug
% form is per publisher -- do not pattern-match one book's from the other's.)
\SetDefaultText{openstax}


%%%%%%%%%%%%%%%%%%%%%%%%%%%%%%%%%%%%%%%%
% 4. SECTIONS CROSS-REFERENCED IN THIS TEXT
%%%%%%%%%%%%%%%%%%%%%%%%%%%%%%%%%%%%%%%%
% LIVE ENTRIES GO HERE, and there are none yet. A lecture's OWN section needs no
% line -- it is declared in the lecture, with \booksection. These are only for
% \textbookref, pointing at some OTHER section.
%
% THESE BELONG TO THE TEXT, NOT TO THE COURSE, which is the easy thing to get
% wrong. A Module A lecture never loads apex.tex, so every OpenStax section a
% Module A lecture cites has to be declared HERE, even if apex.tex's commented
% table happens to list something with the same number.
%
% To add one, copy a line up from the reference block at the foot of this file,
% DELETING THE LEADING "%   ". One entry per line, so these stay greppable
% without compiling anything.
%
% DECLARE ONLY WHAT YOU CITE. The sparseness is a CHECK: \textbookref reports a
% key it cannot resolve and prints the reference unlinked, so a typo is
% something you SEE. Declare the whole book and that same typo silently emits a
% perfect-looking link to the wrong section.
%
% THE GAP THIS FILE USED TO RECORD IS CLOSED. It said that TOPICS.tsv has APEX
% lectures citing OpenStax sections (B.3 cites O1.2, B.4 cites O1.1) and every
% Module A row citing the OpenStax Algebra & Trigonometry text (the [xN.M]
% refs), and that \textbookref could express neither, so a cross-book pointer
% had to be a hand-typed \href in the lecture's prepwork box.
%
% \textbookref now takes the book: \textbookref[algtrig]{3.2}{\S3.2}. Sections
% of a book NO document is built on are declared in books.tex, beside the
% \DeclareBook they qualify; sections of THIS book stay here, unnamed. The
% Module A prepwork boxes were converted when that landed.
%
% Going the other way -- an APEX lecture citing an OpenStax section -- add a
% \DeclareSection[openstax] line to books.tex when a lecture actually cites one.
% It does not belong here, because the document that needs it never loads this
% file.


% ---------------------------------------------------------------------------
% REFERENCE -- MODULE A'S SECTIONS, DELIBERATELY COMMENTED OUT
% ---------------------------------------------------------------------------
% NOTHING BELOW IS LIVE. It is data: the copy-paste source for §4 above, and
% what each Module A lecture's \booksection link= value was taken from.
%
% Only chapter 1 is listed. This course teaches nothing else from OpenStax, so
% unlike apex.tex's whole-book table there is no reason to carry the rest.
%
% !! 1.3 IS VERIFIED -- it was lifted from a live \href in last year's
% !! os1-3content.tex. The other four are pattern-derived from OpenStax's slug
% !! convention (section number, then the hyphenated section title) and are
% !! UNVERIFIED against the live site. A wrong slug still produces a perfectly
% !! good-looking link to a 404, so check one by hand before leaning on it in
% !! front of a class.
%
% --- Chapter 1: Functions and Graphs ---
%   \DeclareSection{1.1}{pages/1-1-review-of-functions}                      % Review of Functions
%   \DeclareSection{1.2}{pages/1-2-basic-classes-of-functions}               % Basic Classes of Functions
%   \DeclareSection{1.3}{pages/1-3-trigonometric-functions}                  % Trigonometric Functions   [VERIFIED]
%   \DeclareSection{1.4}{pages/1-4-inverse-functions}                        % Inverse Functions
%   \DeclareSection{1.5}{pages/1-5-exponential-and-logarithmic-functions}    % Exponential and Logarithmic Functions
